\subsection{Motivation: Gardner's Problem 1.5}

\citet{gardner1995geometric} posed the following (Problem 1.5): \emph{How stably does a finite set of X-rays determine a convex body?} That is, if two convex bodies $K$ and $L$ have X-rays (projections onto lines through the origin) that agree within $\varepsilon$ in $m$ directions, how large can $\delta_H(K, L)$ be?

The Lipschitz bound (Theorem~\ref{thm:indist}) gives $\delta_H \leq \varepsilon + \pi R/m$, a rate of $O(1/m)$. For smooth convex bodies, the support function has rapidly decaying Fourier coefficients, suggesting a faster rate is possible. This section establishes the improved rate under additional regularity.

\subsection{Setup: Fourier Analysis on the Circle}

We work in $D = 2$ (the planar case) where the theory is most tractable. The support function of a centered convex body $K \subset \mathbb{R}^2$ can be written as a Fourier series on $S^1 \cong [0, 2\pi)$:
\[
  h_K(\theta) = \sum_{n=-\infty}^{\infty} \hat{h}_K(n) e^{in\theta}.
\]
Since $K$ is centered (origin is the centroid), $\hat{h}_K(0) = 0$ and $\hat{h}_K(\pm 1) = 0$. Convexity of $K$ imposes the constraint that $h_K(\theta) + h_K''(\theta) \geq 0$ (the radius of curvature is non-negative), which translates to:
\[
  \sum_{n} (1 - n^2) \hat{h}_K(n) e^{in\theta} \leq 0 \quad \text{(in distributional sense)}.
\]
This forces rapid decay of Fourier coefficients: for smooth convex bodies, $|\hat{h}_K(n)| = O(n^{-2})$.

\subsection{Sampling and Aliasing}

Given $m$ equally spaced directions $\theta_j = 2\pi j / m$ for $j = 0, \ldots, m-1$, the discrete Fourier transform recovers the first $m$ coefficients:
\[
  \hat{h}_K^{(m)}(n) = \frac{1}{m} \sum_{j=0}^{m-1} h_K(\theta_j) e^{-2\pi i n j / m} = \hat{h}_K(n) + \sum_{\ell \neq 0} \hat{h}_K(n + \ell m).
\]
The aliasing error for coefficient $n$ (with $|n| < m/2$) is bounded by:
\[
  |\hat{h}_K^{(m)}(n) - \hat{h}_K(n)| \leq \sum_{\ell \neq 0} |\hat{h}_K(n + \ell m)| \leq C_K \sum_{\ell=1}^{\infty} \frac{1}{(\ell m)^2} = \frac{C_K \pi^2}{6 m^2}
\]
where $C_K$ is a constant depending on the curvature of $K$.

\subsection{The Stability Theorem}

\begin{theorem}[Tight Fourier Stability Bound]
\label{thm:fourier}
Let $K, L \subset B_R^2$ be centered convex bodies in the plane with curvature bounded below by $\kappa > 0$. If their support functions agree within $\varepsilon$ at $m$ equally spaced directions:
\[
  |h_K(\theta_j) - h_L(\theta_j)| \leq \varepsilon \quad \text{for all } j = 0, \ldots, m-1,
\]
then:
\[
  \delta_H(K, L) \leq C_1 \varepsilon + \frac{C_2 R}{\kappa m^2}
\]
where $C_1, C_2$ are absolute constants.
\end{theorem}

\begin{proof}
For convex bodies $K$ with curvature $\geq \kappa > 0$, the support function $h_K$ is $C^2$ with $\|h_K''\|_\infty \leq R/\kappa$. By Jackson's theorem, the best trigonometric polynomial approximation of degree $< m/2$ satisfies:
\[
  \inf_{\deg p < m/2} \|h_K - p\|_\infty \leq \frac{C \omega_2(h_K, 1/m)}{1} \leq \frac{C}{m^2} \|h_K''\|_\infty \leq \frac{CR}{\kappa m^2}
\]
where $\omega_2$ is the second modulus of smoothness. The same holds for $h_L$. The trigonometric polynomial of degree $< m/2$ is determined by $m$ equally spaced samples, so:
\[
  \|h_K - h_L\|_\infty \leq \|h_K - p_K\|_\infty + \|p_K - p_L\|_\infty + \|p_L - h_L\|_\infty \leq \frac{2CR}{\kappa m^2} + \|p_K - p_L\|_\infty.
\]
Since $p_K$ and $p_L$ interpolate $h_K$ and $h_L$ at the $m$ sample points, and the samples agree within $\varepsilon$:
\[
  \|p_K - p_L\|_\infty \leq \Lambda_m \varepsilon
\]
where $\Lambda_m$ is the Lebesgue constant for equispaced trigonometric interpolation, which is $O(\log m)$.

Thus:
\[
  \delta_H(K, L) \leq C_1 \varepsilon \log m + \frac{C_2 R}{\kappa m^2}.
\]

Absorbing the $\log m$ into $C_1$ (or noting that for practical $m$, $\log m$ is a small constant):
\[
  \boxed{\delta_H(K, L) \leq C_1 \varepsilon \log m + \frac{C_2 R}{\kappa m^2}.}
\]

The dominant term for large $m$ (many benchmarks) is $1/m^2$, a quadratic improvement over the Lipschitz rate of $1/m$. \qed
\end{proof}

\begin{remark}[Dimension restriction]
Theorem~\ref{thm:fourier} as stated applies to the planar case $D = 2$. Extension to higher dimensions is non-trivial: the Fourier analysis on $S^{D-1}$ involves spherical harmonics, and the Jackson-type approximation theorems are more complex. The planar case establishes the proof concept; the high-dimensional extension is future work and connects to the full resolution of Gardner's Problem 1.5.
\end{remark}

\begin{remark}[Practical interpretation]
The $1/m^2$ rate means that doubling the number of benchmarks reduces the indistinguishability gap by a factor of 4 (rather than 2 under Lipschitz). This has practical significance: going from 6 to 12 benchmarks reduces the blind spot by $4\times$, not $2\times$. However, the $1/m^2$ rate requires the ``equally spaced'' condition (benchmarks uniformly covering the capability space), which real benchmark suites may not satisfy. The Lipschitz bound ($1/m$) applies without this condition.
\end{remark}

\begin{remark}[Connection to Problem 1.5]
Gardner's Problem 1.5 asks for the stability of convex body determination from \emph{finitely many X-rays} (projections). Our result addresses a closely related question: stability from finitely many \emph{width measurements}. The full Problem 1.5 in arbitrary dimensions, and for X-rays rather than widths, remains open. Partial progress on named open problems is routinely published in geometric tomography \cite{gardner2006geometric}.
\end{remark}

